\documentclass[11pt,a4paper]{article}
\usepackage[T1]{fontenc}
\usepackage[margin=26mm,headheight=15pt]{geometry}
\usepackage{amsmath,amsthm,mathtools}
\usepackage{newtxtext,newtxmath}
\usepackage{microtype,booktabs,array,longtable}
\usepackage{enumitem,fancyhdr,xurl}
\usepackage[hidelinks]{hyperref}
\usepackage[nameinlink,noabbrev]{cleveref}
\hypersetup{pdftitle={Vertical-Line Hurwitz Pairings and Exact Stieltjes--Gamma Calculus},pdfauthor={ProveIt research continuation},pdfsubject={Exact identities, rational vertical dilations, compensated Stieltjes functions, and Gamma primitives}}
\pagestyle{fancy}\fancyhf{}
\fancyhead[L]{\small Vertical-line Hurwitz pairings}
\fancyhead[R]{\small ProveIt research continuation}
\fancyfoot[C]{\thepage}
\renewcommand{\headrulewidth}{0.3pt}
\newtheorem{theorem}{Theorem}[section]
\newtheorem{lemma}[theorem]{Lemma}
\newtheorem{proposition}[theorem]{Proposition}
\newtheorem{corollary}[theorem]{Corollary}
\theoremstyle{definition}
\newtheorem{definition}[theorem]{Definition}
\newtheorem{remark}[theorem]{Remark}
\newtheorem{question}{Research question}
\numberwithin{equation}{section}
\newcommand{\C}{\mathbb C}
\newcommand{\R}{\mathbb R}
\newcommand{\N}{\mathbb N}
\newcommand{\Rea}{\operatorname{Re}}
\newcommand{\Li}{\operatorname{Li}}
\newcommand{\CT}{\operatorname{CT}}
\newcommand{\Res}{\operatorname{Res}}
\newcommand{\dd}{\,\mathrm d}
\newcommand{\ii}{\mathrm i}
\newcommand{\bpair}[2]{\mathcal V\!\left[#1,#2\right]}
\newcommand{\coeff}[2]{\left[#1\right]#2}
\newcommand{\sourcepath}[1]{\texttt{\detokenize{#1}}}
\newcommand{\pkg}{\texttt{ProveIt\_Vertical\_Stieltjes\_Pairings\_2026-10-10}}
\allowdisplaybreaks[2]
\setlength{\emergencystretch}{2em}
\setlist[enumerate]{leftmargin=*,itemsep=4pt,topsep=5pt}
\setcounter{tocdepth}{1}
\title{\vspace{-12mm}\textbf{Vertical-Line Hurwitz Pairings\\and Exact Stieltjes--Gamma Calculus}\\[3mm]
\large Rational dilations, Bernoulli compensation,\\polylogarithm identities, and normalized primitives}
\author{Research continuation for Vladimir Reshetnikov's ProveIt project}
\date{10 October 2026 \quad\textbullet\quad Version 1.0}
\begin{document}
\maketitle
\begin{abstract}
We study ordinary, absolutely convergent bilinear integrals in which the
\emph{argument} of a Hurwitz zeta function moves along a vertical line.
Subtracting its spatial principal term produces a Laplace kernel whose
products close in a finite alphabet of Hurwitz zeta functions. Differentiation
at the spectral pole yields an arbitrary-index pairing formula for the
compensated Stieltjes functions
$\gamma_m(z)+\log^{m+1}z/(m+1)$. A finite Bernoulli-subtraction calculus extends
this closure to negative spectral orders and to normalized antiderivatives
of the Stirling remainder. Rationally related vertical scales admit a finite
residue-grid reduction; we give an explicit all-order coefficient formula,
not merely a meromorphic prescription. Additional identities couple these
families to finite generalized harmonic sums, polygamma functions, Lerch
transcendents, and polylogarithm order derivatives. Three illustrative
specializations evaluate a compensated digamma norm, a Stirling-remainder
norm, and a two-scale pairing containing $\log\Gamma(1/4)$.
All general identities are proved by Laplace--Plancherel, finite kernel
algebra, and holomorphic continuation. Reproducible symbolic and numerical
checks are separated from those proofs. A concrete nonreal-argument failure
of the installed mpmath 1.3.0 Stieltjes routine is documented and avoided by
an independent Cauchy-coefficient calculation. No claim of an exhaustive
literature audit, proof-assistant verification, or resolution of the
repository's $S_6$ and revised $S_8$ conjectures is made.
\end{abstract}
\tableofcontents
\clearpage
\section{Scope, relation to the repository, and principal results}
\label{sec:scope}

The canonical manuscript in \cite{repo} already develops a substantial
periodic Hurwitz--Stieltjes convolution calculus. In particular, its
Fourier-normalized family on $\R/\mathbb Z$ leads to shifted correlations,
finite parts, coincident limits, contact terms, and primitive ladders.
Repeating those formulas under new notation would not advance that program.
Here the integration variable is instead the imaginary part of the
\emph{second} Hurwitz argument:
\begin{equation}\label{eq:vertical-scope}
 \frac1{2\pi}\int_{-\infty}^{\infty}
       f(a+\ii y)g(b-\ii y)\dd y,
 \qquad \Rea a,\Rea b>0.
\end{equation}
This is neither an integral in the spectral variable $s=\sigma+\ii y$ nor
a periodic correlation on $[0,1]$. It requires a different subtraction,
a different transform, and different convergence hypotheses. The auxiliary
Hurwitz products over $[0,1]$ studied in \cite{shpot2016,shpot-paris} also
have a different integration geometry and use a different auxiliary
function. Our starting transform identities are classical; the contribution
of this note is the explicit, jointly proved closure calculus and its
specializations. We do not assert priority for every displayed identity.

Set
\begin{equation}\label{eq:Z-definition-intro}
 Z_s(z)=\zeta(s,z)-\frac{z^{1-s}}{s-1},
 \qquad
 G_m(z)=\gamma_m(z)+\frac{\log^{m+1}z}{m+1}.
\end{equation}
The value at $s=1$ is removable, and $G_0(z)=\log z-\psi(z)$.
The main equal-scale theorem evaluates \eqref{eq:vertical-scope} for
$f=Z_s$, $g=Z_t$ by a one-variable function $K(s+t-1,a+b)$.
Its all-order Taylor coefficients then evaluate every pairing of $G_m$
and $G_n$. Finite Bernoulli subtractions give the same conclusion for
negative spectral orders. Rational vertical dilations introduce only a
finite set of shifted Hurwitz arguments, rather than an irreducible double
zeta object.

Three consequences indicate the scale and normalization of the results.
Define the analytic Stirling remainder in the right half-plane by
\begin{equation}\label{eq:Omega-intro}
 \Omega(z)=\log\Gamma(z)-(z-\tfrac12)\log z+z-\tfrac12\log(2\pi).
\end{equation}
Then
\begin{align}
 \frac1{2\pi}\int_{\R}
   |G_0(\tfrac12+\ii y)|^2\dd y
   &=\log(2\pi)-\frac32,\label{eq:headline-digamma}\\
 \frac1{2\pi}\int_{\R}
   |\Omega(\tfrac12+\ii y)|^2\dd y
   &=\zeta'(-1)+\frac{\zeta(3)}{2\pi^2}+\frac19,
       \label{eq:headline-stirling}\\
 \frac1{2\pi}\int_{\R}
    G_0(\tfrac12+\ii y)G_0(\tfrac12-2\ii y)\dd y
   &=\log\Gamma(\tfrac14)-\tfrac12\log\pi-\tfrac14\log2
       \notag\\[-1mm]
   &\hspace{9mm}+\tfrac34\log3-\tfrac34-\tfrac\pi8.
       \label{eq:headline-quarter}
\end{align}
Their right-hand sides are approximately $0.33787706640934548$,
$0.0065868815274467235$, and $0.22363092143574916$, respectively.
These numerical values are diagnostics; the proofs appear below.

\subsection{What is and is not resolved}
The questions answered here are precise: whether spatially compensated
vertical Stieltjes pairings close at arbitrary indices; whether negative
spectral orders can be handled without assigning finite parts to the
vertical integral; and whether rationally related slopes preserve finite
Hurwitz closure. The answer to each is affirmative, with explicit formulas
and proofs. These are newly developed continuations in this package, not a
claim that a named conjecture in an unread incoming archive has been solved.
The repository's $S_6$ and revised $S_8$ problems are not settled here.

The audit was targeted. The canonical README, the periodic convolution
foundation, a terminal rational-reduction section, the incoming README,
and incoming metadata were inspected. Not every recent ZIP was decoded or
audited. Thus absence of a formula from these inspected sections is not a
proof of absence from the entire corpus. The package includes an integration
map and a source inventory so that an independent maintainer can perform
that final overlap check.

\subsection{Conventions that affect the answer}
Throughout, $\log z$ is the principal logarithm on $\Rea z>0$.
The function $\log\Gamma(z)$ means its analytic logarithm, real on the
positive real axis; it does \emph{not} mean the pointwise principal value
of $\log(\Gamma(z))$. A prime on $\zeta$ or $\Phi$ denotes differentiation
in the spectral order, unless a different variable is explicitly indicated.
We use
\begin{equation}\label{eq:laurent-convention}
 \zeta(1+u,z)=\frac1u+
      \sum_{m\ge0}\frac{(-1)^m}{m!}\gamma_m(z)u^m.
\end{equation}
All the vertical integrals asserted as theorems are ordinary absolutely
convergent Lebesgue integrals. The words ``constant term'' and ``removable
value'' used later refer to a finite \emph{spectral expression}, not to a
regularization of a divergent integral in $y$.

\section{Laplace--Plancherel and the basic Hurwitz closure}
\label{sec:laplace}

\subsection{A bilinear transform lemma}
For a function supported on $(0,\infty)$ use the Fourier convention
$\widehat f(y)=\int_0^\infty f(x)e^{-\ii yx}\dd x$.
The following version of Plancherel fixes every scale factor needed later.

\begin{lemma}[Bilinear Laplace--Plancherel]\label{thm:plancherel}
Suppose $f,g\in L^1(0,\infty)\cap L^2(0,\infty)$. Then
\begin{equation}\label{eq:bilinear-plancherel}
 \frac1{2\pi}\int_{\R}\widehat f(y)\widehat g(-y)\dd y
     =\int_0^\infty f(x)g(x)\dd x.
\end{equation}
Both products are absolutely integrable. For $p,q>0$,
\begin{equation}\label{eq:scaled-plancherel}
 \frac1{2\pi}\int_{\R}\widehat f(py)\widehat g(-qy)\dd y
     =\int_0^\infty f(qx)g(px)\dd x.
\end{equation}
If $f$ and $g$ depend holomorphically on parameters as $L^2$ vectors, all
parameter derivatives can be taken inside the bilinear pairing.
\end{lemma}
\begin{proof}
Apply the usual Hermitian Plancherel identity to $f$ and $\overline g$,
noting that $\overline{\widehat{\overline g}(y)}=\widehat g(-y)$.
Cauchy--Schwarz gives absolute integrability. For the scaled formula,
$\widehat f(py)$ is the Fourier transform of $p^{-1}f(x/p)$.
Equation~\eqref{eq:bilinear-plancherel} therefore gives
$(pq)^{-1}\int_0^\infty f(x/p)g(x/q)\dd x$; substitute $x=pqu$.
Finally, the bilinear map $(f,g)\mapsto\int fg$ is continuous on
$L^2\times L^2$. Its composition with holomorphic vector-valued maps is
holomorphic and satisfies the usual product rule.
\end{proof}

\subsection{Spatial compensation}
Write
\begin{equation}\label{eq:h-basic}
 Q(x)=\frac1{1-e^{-x}},\qquad h(x)=Q(x)-\frac1x.
\end{equation}
At the origin, $h(x)=\frac12+\frac{x}{12}+O(x^3)$; at infinity it has at
most polynomial growth. The standard Mellin representation of the Hurwitz
zeta function \cite{dlmf-zeta}, after subtracting the elementary Laplace
transform of $x^{s-2}$, gives
\begin{equation}\label{eq:Z-laplace}
 Z_s(z)=\frac1{\Gamma(s)}\int_0^\infty
            x^{s-1}e^{-zx}h(x)\dd x,
 \qquad \Rea s>0,\quad \Rea z>0.
\end{equation}
Initially one can subtract the two integrals for $\Rea s>1$.
The combined integrand is integrable for $\Rea s>0$ and is locally
uniformly dominated there, which proves the continuation. In particular,
the normalized subtraction is spatial: removing only $1/(s-1)$ would
not produce \eqref{eq:Z-laplace}.

For $\Rea s>1/2$, the function $x^{s-1}e^{-ax}h(x)$ lies in
$L^1\cap L^2$. On each compact subset of this parameter region, every
fixed spectral derivative merely adds a power of $\log x$, which is still
locally dominated by an integrable power at zero and an exponential at
infinity. This verifies the holomorphic-vector hypothesis of
\cref{thm:plancherel}; it will justify all subsequent Stieltjes jet
extractions without a separate unjustified interchange at each order.

\begin{theorem}[Equal-scale compensated Hurwitz pairing]\label{thm:basic}
Let $\Rea a,\Rea b>0$ and $\Rea s,\Rea t>1/2$. Put
$A=a+b$ and $w=s+t-1$. Then
\begin{equation}\label{eq:basic-pair}
 \frac1{2\pi}\int_{\R}Z_s(a+\ii y)Z_t(b-\ii y)\dd y
   =\frac{\Gamma(w)}{\Gamma(s)\Gamma(t)}K(w,A),
\end{equation}
where
\begin{equation}\label{eq:K}
 \boxed{\displaystyle
 K(w,A)=\left(1-\frac2{w-1}\right)\zeta(w-1,A)
 +(1-A)\zeta(w,A)
 +\frac{A^{2-w}}{(w-1)(w-2)}.}
\end{equation}
The apparent singularities at $w=1,2$ are removable. The combined function
$\Gamma(w)K(w,A)$ is holomorphic for $\Rea w>0$ and $\Rea A>0$.
\end{theorem}
\begin{proof}
By \eqref{eq:Z-laplace} and \cref{thm:plancherel}, the integral equals
\begin{equation}\label{eq:basic-moment}
 \frac1{\Gamma(s)\Gamma(t)}
   \int_0^\infty x^{w-1}e^{-Ax}h(x)^2\dd x.
\end{equation}
For $\Rea w>2$, expand $h^2=Q^2-2Q/x+x^{-2}$.
The exact power series
\[
 Q(x)^2=\sum_{n\ge0}(n+1)e^{-nx}
\]
yields
\[
 \int_0^\infty x^{w-1}e^{-Ax}Q(x)^2\dd x
 =\Gamma(w)\{\zeta(w-1,A)+(1-A)\zeta(w,A)\},
\]
because $n+1=(n+A)+(1-A)$. The other two integrals are
$-2\Gamma(w-1)\zeta(w-1,A)$ and $\Gamma(w-2)A^{2-w}$.
Using the Gamma recurrence gives \eqref{eq:K}.
The combined integral \eqref{eq:basic-moment} is holomorphic for
$\Rea w>0$, since $h^2=1/4+O(x)$ at zero. The identity theorem extends
the finite expression to that half-plane and removes its apparent poles.
\end{proof}

The pole cancellation can also be checked directly. At $w=1$ the residue
is
\[
 -2\zeta(0,A)+(1-A)-A=0,
 \qquad \zeta(0,A)=\tfrac12-A.
\]
At $w=2$, the residue $-1$ from the first zeta term cancels the residue
$+1$ from the elementary term. Neither cancellation is a license to
substitute separately into the singular summands.

\begin{corollary}[An exact compensated digamma integral]\label{thm:digamma}
For $\Rea a,\Rea b>0$ and $A=a+b$,
\begin{align}
 &\frac1{2\pi}\int_{\R}
 [\log(a+\ii y)-\psi(a+\ii y)]
 [\log(b-\ii y)-\psi(b-\ii y)]\dd y \notag\\
 &\hspace{8mm}=
 \frac12-2A+A\log A-2\log\Gamma(A)+\log(2\pi)
       -(1-A)\psi(A).\label{eq:digamma-pair}
\end{align}
\end{corollary}
\begin{proof}
The constant term of \eqref{eq:K} at $w=1$ is
\[
 \zeta(0,A)-2\zeta'(0,A)+(1-A)\gamma_0(A)+A(\log A-1).
\]
Use $\gamma_0=-\psi$ and
$\zeta'(0,A)=\log\Gamma(A)-\frac12\log(2\pi)$
\cite{dlmf-zeta}. Formula~\eqref{eq:headline-digamma} follows at $A=1$.
\end{proof}

\subsection{Position derivatives}
Direct differentiation gives the exact ladder
\begin{equation}\label{eq:Z-position-ladder}
 \partial_z Z_s(z)=-sZ_{s+1}(z),\qquad
 \partial_z^r Z_s(z)=(-1)^r(s)_r Z_{s+r}(z).
\end{equation}
Consequently, for integers $r,k\ge0$, differentiating
\eqref{eq:basic-pair} in $a$ and $b$ gives
\begin{align}
 &\frac1{2\pi}\int_{\R}
    (\partial_z^r Z_s)(a+\ii y)(\partial_z^k Z_t)(b-\ii y)\dd y
       \notag\\
 &\quad=\partial_A^{r+k}
       \frac{\Gamma(w)K(w,A)}{\Gamma(s)\Gamma(t)}
   =\frac{(-1)^{r+k}\Gamma(w+r+k)}{\Gamma(s)\Gamma(t)}
                K(w+r+k,A).\label{eq:position-pair}
\end{align}
The last equality follows by differentiating the Laplace moment, and is
valid at all its removable spectral values. Thus the differentiation
calculus closes without introducing a second independent position
parameter. Spectral differentiation of this identity gives the same result
for derivatives of arbitrary $G_m$.

\section{Arbitrary-index Stieltjes pairings}
\label{sec:stieltjes}

\subsection{The complete coefficient theorem}
Expansion of the subtracted term in \eqref{eq:Z-definition-intro} gives
\begin{equation}\label{eq:Z-G-expansion}
 Z_{1+u}(z)=\sum_{m\ge0}\frac{(-1)^m}{m!}G_m(z)u^m.
\end{equation}
For $\Rea a,\Rea b>0$, define
\begin{equation}\label{eq:Imn}
 I_{m,n}(A)=\frac1{2\pi}\int_{\R}
            G_m(a+\ii y)G_n(b-\ii y)\dd y,
 \qquad A=a+b.
\end{equation}
The following theorem both proves that this depends only on $A$ and
specifies a finite answer at every pair of indices.

\begin{theorem}[All-index Stieltjes closure]\label{thm:all-stieltjes}
For all integers $m,n\ge0$, the integral \eqref{eq:Imn} converges absolutely
and
\begin{equation}\label{eq:all-jets}
 \boxed{\displaystyle
 I_{m,n}(A)=(-1)^{m+n}m!n!\,[u^mv^n]
      C(u,v)K(1+u+v,A),}
\end{equation}
where
\begin{align}
 C(u,v)&=\frac{\Gamma(1+u+v)}{\Gamma(1+u)\Gamma(1+v)},\label{eq:C}\\
 \log C(u,v)&=\sum_{r\ge2}\frac{(-1)^r\zeta(r)}r
          \{(u+v)^r-u^r-v^r\}.\label{eq:log-C}
\end{align}
Write $K(1+\epsilon,A)=\sum_{\ell\ge0}k_\ell(A)\epsilon^\ell$.
Its coefficients are
\begin{align}
 k_\ell(A)={}&\frac{\zeta^{(\ell)}(0,A)}{\ell!}
  -\frac{2\zeta^{(\ell+1)}(0,A)}{(\ell+1)!}
  +(1-A)\frac{(-1)^\ell\gamma_\ell(A)}{\ell!}\notag\\
 &-A\sum_{j=0}^{\ell+1}\frac{(-\log A)^j}{j!}.
       \label{eq:k-coefficients}
\end{align}
In particular, $I_{m,n}$ belongs to the finite algebra generated by
$A$, $\log A$, the constants $\zeta(2),\ldots,\zeta(m+n)$,
$\gamma_\ell(A)$ for $\ell\le m+n$, and
$\zeta^{(j)}(0,A)$ for $j\le m+n+1$.
\end{theorem}
\begin{proof}
The map $u\mapsto Z_{1+u}(a+\ii y)$ is holomorphic as an $L^2(\R)$ vector
in a neighborhood of zero, by the estimates preceding \cref{thm:basic}.
Its derivatives are precisely $(-1)^mG_m$. Hence every $G_m$ belongs to
$L^2$, and the product in \eqref{eq:Imn} is absolutely integrable.
Taking the two Taylor coefficients in \eqref{eq:basic-pair} gives
\eqref{eq:all-jets}. The Taylor expansion of $\log\Gamma(1+u)$ proves
\eqref{eq:log-C}; its linear Euler-constant terms cancel.

For the explicit coefficients, write
\[
 K(1+\epsilon,A)=\left(1-\frac2\epsilon\right)\zeta(\epsilon,A)
 +(1-A)\zeta(1+\epsilon,A)
 +\frac{A e^{-\epsilon\log A}}{\epsilon(\epsilon-1)}.
\]
The pole coefficient vanishes as already checked. Since
\[
 \frac{Ae^{-\epsilon\log A}}{\epsilon(\epsilon-1)}
 =-A\sum_{r\ge0}\epsilon^{r-1}
          \sum_{j=0}^r\frac{(-\log A)^j}{j!},
\]
its coefficient of $\epsilon^\ell$ is the last term of
\eqref{eq:k-coefficients}. The other terms follow from the Taylor and
Laurent series of the two zeta factors. Only total degree at most $m+n$
of $C$ is used, proving the finite-alphabet assertion.
\end{proof}

This theorem is an actual finite recipe. If
$C(u,v)=\sum c_{r,k}u^rv^k$, then equivalently
\begin{equation}\label{eq:finite-jet-recipe}
 I_{m,n}(A)=(-1)^{m+n}m!n!
  \sum_{r=0}^m\sum_{k=0}^n c_{r,k}
       \binom{m+n-r-k}{m-r}k_{m+n-r-k}(A).
\end{equation}
No two-variable special function or undetermined integration constant
remains. This is a closure statement, not a claim that all its generators
are algebraically independent or reducible to elementary constants.

\subsection{Low-order identities and a useful degeneration}
Let $\mu_j(A)=\partial_w^jK(1,A)=j!k_j(A)$. Direct coefficient extraction
from \eqref{eq:log-C} gives
\begin{align}
 I_{0,0}&=\mu_0,&
 I_{m,0}&=(-1)^m\mu_m,\label{eq:jet-row0}\\
 I_{1,1}&=\mu_2+\zeta(2)\mu_0,\label{eq:jet11}\\
 I_{2,1}&=-\mu_3-2\zeta(2)\mu_1+2\zeta(3)\mu_0,\label{eq:jet21}\\
 I_{3,1}&=\mu_4+3\zeta(2)\mu_2-6\zeta(3)\mu_1
                   +6\zeta(4)\mu_0,\label{eq:jet31}\\
 I_{2,2}&=\mu_4+4\zeta(2)\mu_2-8\zeta(3)\mu_1
                   +\{6\zeta(4)+2\zeta(2)^2\}\mu_0.
                   \label{eq:jet22}
\end{align}
Here and below, every $\mu_j$ in one equation has the same argument $A$.
For example, the coefficient of $u^2v^2$ in $C$ is
$\frac32\zeta(4)+\frac12\zeta(2)^2$; the factor $2!2!$ explains the last
coefficient in \eqref{eq:jet22}. This factorial is easy to lose if ordinary
and exponential generating functions are mixed.

At $A=1$, the entire explicit Stieltjes term in
\eqref{eq:k-coefficients} disappears:
\begin{equation}\label{eq:A-one-degeneration}
 k_\ell(1)=\frac{\zeta^{(\ell)}(0)}{\ell!}
          -\frac{2\zeta^{(\ell+1)}(0)}{(\ell+1)!}-1.
\end{equation}
Thus every vertical pairing with $a+b=1$ is expressible using Riemann zeta
order derivatives at zero and ordinary positive zeta values. This does
not say that the individual Stieltjes functions disappear pointwise;
the simplification occurs after pairing.

\subsection{A second, kernel-polynomial description}
Define real-coefficient polynomials $E_m(L)$ by
\begin{equation}\label{eq:E-generating}
 \sum_{m\ge0}E_m(L)\frac{v^m}{m!}
          =\frac{e^{-vL}}{\Gamma(1-v)}.
\end{equation}
For example,
\[
 E_0=1,\qquad E_1=-(L+\gamma),\qquad
 E_2=(L+\gamma)^2-\zeta(2).
\]
Differentiating \eqref{eq:Z-laplace} gives
\begin{align}
 G_m(z)&=\int_0^\infty e^{-zx}h(x)E_m(\log x)\dd x,
               \label{eq:G-bell-laplace}\\
 I_{m,n}(A)&=\int_0^\infty e^{-Ax}h(x)^2
                  E_m(\log x)E_n(\log x)\dd x.
               \label{eq:Stieltjes-bell-pair}
\end{align}
All these differentiations are justified by the same compact-parameter
bounds used in the proof of \cref{thm:all-stieltjes}.
Equations~\eqref{eq:all-jets} and \eqref{eq:Stieltjes-bell-pair} provide two
useful computational routes: the former uses a finite zeta/Stieltjes
alphabet, while the latter uses one rapidly decaying quadrature with a kernel regular for positive $x$. They are independently implemented in the verification code.

The equal-scale identity also supplies all derivatives in the horizontal
arguments: one differentiates \eqref{eq:all-jets} in $A$, or inserts powers
of $-x$ in \eqref{eq:Stieltjes-bell-pair}. No new spectral-order restriction
arises from a nonnegative number of such derivatives.

\subsection{A complex-safe Hermite identity}
The generalized Stieltjes parameter must be treated holomorphically, not by
applying a real-part formula outside its real domain. A useful exact
representation, valid at every index, is the following.

\begin{proposition}[Holomorphic Hermite representation]\label{thm:complex-Hermite}
For $\Rea z>0$ and every integer $m\ge0$,
\begin{align}
 G_m(z)={}&\frac{\log^m z}{2z}
  +\ii\int_0^\infty\frac1{e^{2\pi t}-1}
       \left\{\frac{\log^m(z+\ii t)}{z+\ii t}
                  -\frac{\log^m(z-\ii t)}{z-\ii t}\right\}\dd t.
                     \label{eq:complex-Hermite}
\end{align}
The integral is absolutely convergent and holomorphic in $z$ throughout
the right half-plane. For nonreal $z$, neither term may be replaced by
the complex conjugate of the other.
\end{proposition}
\begin{proof}
The standard Hermite representation of Hurwitz zeta
\cite{dlmf-zeta}, after subtracting its spatial principal term, is
\[
 Z_s(z)=\frac1{2z^s}
   +\ii\int_0^\infty
        \frac{(z+\ii t)^{-s}-(z-\ii t)^{-s}}{e^{2\pi t}-1}\dd t.
\]
It can equivalently be obtained from \eqref{eq:Z-laplace} using
$h(x)=1/2+2\int_0^\infty\sin(tx)/(e^{2\pi t}-1)\dd t$.
The two powers have a difference $O(t)$ at zero, uniformly on compact
subsets of $s\in\C$, $\Rea z>0$; the denominator also has a simple zero.
At infinity the exponential denominator dominates all powers and
logarithms arising from any fixed spectral derivative. Therefore the
combined integral is jointly holomorphic and may be differentiated in
$s$ at $s=1$. Multiplying the $m$th derivative by $(-1)^m$ proves
\eqref{eq:complex-Hermite}. Holomorphic continuation in $z$ from positive
real values gives the entire right half-plane; the paired expression,
not a real-part projection, is what makes this continuation legitimate.
\end{proof}

This representation is classical in origin, but spelling out its paired
complex form is important for the verification of the present vertical-line
identities. \Cref{sec:software-audit} records a concrete numerical failure
when a real-part implementation is used at nonreal $z$, and supplies both
\eqref{eq:complex-Hermite} and Cauchy-coefficient extraction as repairs.

\section{Polygamma functions and generalized harmonic sums}
\label{sec:harmonics}

\subsection{Uncompensated polygamma pairings}
For uncompensated Hurwitz zeta functions, the Laplace kernel $Q(x)$ has a
simple pole at zero. Consequently, $\Rea s,\Rea t>3/2$ is a sufficient
$L^2$ region for the same argument to give
\begin{equation}\label{eq:plain-zeta-pair}
 \frac1{2\pi}\int_{\R}\zeta(s,a+\ii y)\zeta(t,b-\ii y)\dd y
 =\frac{\Gamma(w)}{\Gamma(s)\Gamma(t)}
       \{\zeta(w-1,A)+(1-A)\zeta(w,A)\}.
\end{equation}
This follows directly from the $Q^2$ calculation in the proof of
\cref{thm:basic}; no compensation is needed in this smaller domain.
Using the standard identity
$\psi^{(r)}(z)=(-1)^{r+1}r!\zeta(r+1,z)$ gives the following consequence.

\begin{corollary}\label{thm:polygamma}
For integers $r,k\ge1$, $\Rea a,\Rea b>0$, and $A=a+b$,
\begin{align}
 &\frac1{2\pi}\int_{\R}\psi^{(r)}(a+\ii y)
                                  \psi^{(k)}(b-\ii y)\dd y\notag\\
 &\quad=(-1)^{r+k}(r+k)!
      \{\zeta(r+k,A)+(1-A)\zeta(r+k+1,A)\}.\label{eq:polygamma-pair}
\end{align}
In particular,
\begin{equation}\label{eq:polygamma-norm}
 \frac1{2\pi}\int_{\R}
       |\psi^{(r)}(\tfrac12+\ii y)|^2\dd y=(2r)!\zeta(2r).
\end{equation}
\end{corollary}
The restriction $r,k\ge1$ matters. Setting both to zero would introduce
the divergent squared norm of the uncorrected digamma function, not the
finite value in \eqref{eq:headline-digamma}.

\subsection{Finite generalized harmonic families}
For a positive integer $N$, define
\begin{equation}\label{eq:finite-H}
 \mathcal H_N^{(s)}(z)=\sum_{j=0}^{N-1}(z+j)^{-s}.
\end{equation}
At $z=1$ this is the usual generalized harmonic number $H_N^{(s)}$.
For $\Rea s>0$, it is the Laplace transform with kernel
$\sum_{j=0}^{N-1}e^{-jx}$, divided by $\Gamma(s)$.

\begin{theorem}[Finite harmonic and mixed pairings]\label{thm:harmonic}
Let $M,N\ge1$, $\Rea a,\Rea b>0$, $\Rea s,\Rea t>1/2$, and
$A=a+b$, $w=s+t-1$. Then
\begin{align}
 &\frac1{2\pi}\int_{\R}\mathcal H_M^{(s)}(a+\ii y)
                           \mathcal H_N^{(t)}(b-\ii y)\dd y\notag\\
 &\qquad=\frac{\Gamma(w)}{\Gamma(s)\Gamma(t)}
       \sum_{j=0}^{M-1}\sum_{k=0}^{N-1}(A+j+k)^{-w},\label{eq:HH-pair}\\
 &\frac1{2\pi}\int_{\R}Z_s(a+\ii y)
                           \mathcal H_N^{(t)}(b-\ii y)\dd y\notag\\
 &\qquad=\frac{\Gamma(w)}{\Gamma(s)\Gamma(t)}
                         \sum_{k=0}^{N-1}Z_w(A+k).\label{eq:ZH-pair}
\end{align}
The second expression is taken at its removable value if $w=1$.
Arbitrary spectral and position derivatives are permitted inside both
identities in the stated domain.
\end{theorem}
\begin{proof}
Multiply the two finite exponential kernels and apply
\cref{thm:plancherel}. Each exponential term contributes
$\Gamma(w)(A+j+k)^{-w}$. In the mixed product it contributes
$\Gamma(w)Z_w(A+k)$ by \eqref{eq:Z-laplace}.
The kernels are bounded at zero, so the same $L^2$ holomorphy argument
applies to all derivatives.
\end{proof}

For $a+b=1$ and $s=t=1$, these become particularly elementary:
\begin{align}
 &\frac1{2\pi}\int_{\R}\mathcal H_M^{(1)}(a+\ii y)
                        \mathcal H_N^{(1)}(b-\ii y)\dd y\notag\\
 &\quad=(M+N)H_{M+N}-MH_M-NH_N,\label{eq:HH-elementary}\\
 &\frac1{2\pi}\int_{\R}G_0(a+\ii y)
                        \mathcal H_N^{(1)}(b-\ii y)\dd y\notag\\
 &\quad=\log\Gamma(N+1)-NH_N+N(1+\gamma).
               \label{eq:GH-elementary}
\end{align}
Indeed, the first double sum can be summed along its diagonals, and the
second is $\sum_{k=1}^N(\log k-\psi(k))$.
Taking spectral derivatives supplies finite sums of
$\log^r(A+j+k)/(A+j+k)^w$ and corresponding exact mixed identities for
higher $G_m$, with the same Gamma-quotient coefficient bookkeeping as
in \cref{sec:stieltjes}.

\section{Lerch products and polylogarithm order derivatives}
\label{sec:lerch}

Let $\xi$ lie in $\{\xi:|\xi|\le1,\ \xi\ne1\}$. The Lerch transcendent
has the integral representation \cite{dlmf-lerch}
\begin{equation}\label{eq:lerch-laplace}
 \Phi(\xi,s,z)=\frac1{\Gamma(s)}\int_0^\infty
           \frac{x^{s-1}e^{-zx}}{1-\xi e^{-x}}\dd x,
 \qquad\Rea s>0,\quad\Rea z>0.
\end{equation}
For $|\xi|<1$ this also follows by summing its defining series.
The kernel is bounded at zero and along the positive real axis. Boundary
colors $|\xi|=1$, $\xi\ne1$, follow by dominated continuation from inside
the disk. No assertion here includes the singular color $\xi=1$ without
the separate compensation already developed.

\begin{theorem}[Colored pairings and confluent colors]\label{thm:lerch}
Let $\Rea s,\Rea t>1/2$, $\Rea a,\Rea b>0$, and set
$A=a+b$, $w=s+t-1$. For admissible colors $\xi\ne\eta$,
\begin{align}
 &\frac1{2\pi}\int_{\R}\Phi(\xi,s,a+\ii y)
                             \Phi(\eta,t,b-\ii y)\dd y\notag\\
 &\quad=\frac{\Gamma(w)}{\Gamma(s)\Gamma(t)}
           \frac{\xi\Phi(\xi,w,A)-\eta\Phi(\eta,w,A)}{\xi-\eta}.
               \label{eq:lerch-distinct}
\end{align}
At $\eta=\xi$, the quotient is replaced by
\begin{equation}\label{eq:lerch-confluent}
 \Phi(\xi,w-1,A)+(1-A)\Phi(\xi,w,A).
\end{equation}
In particular, when $A=1$ and $\xi\ne\eta$,
\begin{align}
 &\frac1{2\pi}\int_{\R}\Phi(\xi,s,a+\ii y)
                             \Phi(\eta,t,1-a-\ii y)\dd y\notag\\
 &\quad=\frac{\Gamma(w)}{\Gamma(s)\Gamma(t)}
                 \frac{\Li_w(\xi)-\Li_w(\eta)}{\xi-\eta}.
                     \label{eq:polylog-divdiff}
\end{align}
For equal nonzero colors the final quotient is $\Li_{w-1}(\xi)/\xi$;
at $\xi=0$ its removable value is $1$.
\end{theorem}
\begin{proof}
For distinct colors, use the exact partial fraction identity
\[
 \frac1{(1-\xi r)(1-\eta r)}
 =\frac{\xi/(1-\xi r)-\eta/(1-\eta r)}{\xi-\eta}
\]
in the Laplace product supplied by \cref{thm:plancherel}.
For equal colors, the kernel has series
$\sum_{n\ge0}(n+1)\xi^n e^{-nx}$, and
$n+1=(n+A)+(1-A)$ proves \eqref{eq:lerch-confluent}.
This series argument initially uses $|\xi|<1$; its identity extends to
the stated boundary colors by the same dominated kernel argument.
Finally, $\xi\Phi(\xi,w,1)=\Li_w(\xi)$, including the removable
interpretation at zero. All products converge by $L^2$.
\end{proof}

Order differentiation of \eqref{eq:polylog-divdiff} is a direct source of
identities for $\partial_w^r\Li_w(\xi)$ and mixed order derivatives of
Lerch functions. In particular, at $s=t=1$, the total-order derivatives
are generated by
\[
 C(u,v)\frac{\Li_{1+u+v}(\xi)-\Li_{1+u+v}(\eta)}{\xi-\eta}.
\]
This is a convergent parameter identity near the origin, rather than an
experimental relation inferred from numerical integer-relation searches.

\subsection{Mixed Stieltjes--Lerch closure and the spectral resonance}
\begin{theorem}\label{thm:mixed-lerch}
Under the hypotheses of \cref{thm:lerch},
\begin{align}
 &\frac1{2\pi}\int_{\R}Z_s(a+\ii y)\Phi(\xi,t,b-\ii y)\dd y\notag\\
 &\quad=\frac{\Gamma(w)}{\Gamma(s)\Gamma(t)}
  \left\{\frac{\zeta(w,A)-\xi\Phi(\xi,w,A)}{1-\xi}
                     -\frac{\Phi(\xi,w-1,A)}{w-1}\right\}.
                         \label{eq:mixed-lerch}
\end{align}
The brace is holomorphic for $\Rea w>0$. At $s=t=1$ it gives
\begin{align}
 &\frac1{2\pi}\int_{\R}G_0(a+\ii y)\Phi(\xi,1,b-\ii y)\dd y\notag\\
 &\quad=\frac{-\psi(A)-\xi\Phi(\xi,1,A)}{1-\xi}
                                  -\Phi'(\xi,0,A).
                                      \label{eq:mixed-lerch-resonance}
\end{align}
\end{theorem}
\begin{proof}
The product kernel is
\[
 \frac{h(x)}{1-\xi e^{-x}}
 =\frac{Q(x)-\xi/(1-\xi e^{-x})}{1-\xi}
                     -\frac1{x(1-\xi e^{-x})}.
\]
Integrate the separated expression first for $\Rea w>1$.
The combined kernel is bounded at zero, so its Mellin transform is
holomorphic for $\Rea w>0$. At $w=1$, both singular summands have residue
$1/(1-\xi)$ with opposite signs, because
$\Phi(\xi,0,A)=1/(1-\xi)$. Their constant terms give
\eqref{eq:mixed-lerch-resonance}. Derivatives of this removable identity
are justified by the original holomorphic $L^2$ pairing.
\end{proof}

When $A=1$, \eqref{eq:mixed-lerch-resonance} becomes
\begin{equation}\label{eq:mixed-Li0}
 \frac1{2\pi}\int_{\R}G_0(a+\ii y)\Phi(\xi,1,1-a-\ii y)\dd y
 =\frac{\gamma+\log(1-\xi)}{1-\xi}
       -\frac{\left.\partial_s\Li_s(\xi)\right|_{s=0}}{\xi}.
\end{equation}
At zero color the final quotient is interpreted by continuity. At
$\xi=-1$, the right-hand side reduces to
\begin{equation}\label{eq:alternating-special}
 \frac\gamma2+\log2-\frac12\log\pi.
\end{equation}
Indeed, $\Li_s(-1)=(2^{1-s}-1)\zeta(s)$, whence
$\Li'_0(-1)=\frac12\log(2/\pi)$. Moreover,
\[
 \Phi(-1,1,z)=\frac12\left\{\psi\left(\frac{z+1}{2}\right)
                           -\psi\left(\frac z2\right)\right\}.
\]
Thus \eqref{eq:alternating-special} is also an explicit compensated
digamma--digamma-difference integral.

\subsection{Every nontrivial root of unity gives a finite Gamma formula}
If $\xi^q=1$ and $\xi\ne1$, grouping the Lerch series modulo $q$ gives
\begin{equation}\label{eq:root-unity}
 \Phi(\xi,s,A)=q^{-s}\sum_{r=0}^{q-1}\xi^r
                                \zeta\left(s,\frac{A+r}{q}\right).
\end{equation}
The zeta residues at $s=1$ cancel because $\sum\xi^r=0$. With
$c_r=(A+r)/q$, its two required values are
\begin{align}
 \Phi(\xi,1,A)&=-\frac1q\sum_{r=0}^{q-1}\xi^r\psi(c_r),\label{eq:root-one}\\
 \Phi'(\xi,0,A)&=\sum_{r=0}^{q-1}\xi^r\log\Gamma(c_r)
                                      -\frac{\log q}{1-\xi}.
                                          \label{eq:root-zero-derivative}
\end{align}
Substituting into \eqref{eq:mixed-lerch-resonance} yields the finite family
\begin{align}
 &\frac1{2\pi}\int_{\R}G_0(a+\ii y)\Phi(\xi,1,b-\ii y)\dd y\notag\\
 &\quad=\frac{-\psi(A)+(\xi/q)\sum_{r=0}^{q-1}\xi^r\psi(c_r)
                      +\log q}{1-\xi}
             -\sum_{r=0}^{q-1}\xi^r\log\Gamma(c_r).
                 \label{eq:root-Gamma-pair}
\end{align}
For higher indices, differentiating the finite sum
\eqref{eq:root-unity} gives a triangular combination of Hurwitz order
jets and powers of $\log q$. This is a proved finite reduction at every
index, with no presumption of further transcendence reductions.

\section{Bernoulli compensation and negative spectral orders}
\label{sec:compensation}

\subsection{A hierarchy with an exact convergence domain}
The basic kernel $h$ is bounded but does not vanish at zero. Its Mellin
transform therefore does not directly reach the Gamma logarithm at spectral
order zero. We now subtract exactly the required Taylor terms.
Let $B_{2k}$ be the Bernoulli numbers and put
\begin{align}
 P_0(x)&=x^{-1},\label{eq:P0}\\
 P_d(x)&=x^{-1}+\frac12+
       \sum_{k=1}^{d-1}\frac{B_{2k}}{(2k)!}x^{2k-1},\qquad d\ge1,
                 \label{eq:Pd}\\
 h_d(x)&=Q(x)-P_d(x),\qquad
 \nu_d=\begin{cases}0,&d=0,\\2d-1,&d\ge1.\end{cases}
                 \label{eq:hd}
\end{align}
An empty sum is zero. Thus $h_0=h$, $h_1=h-1/2$, and
$h_2=h-1/2-x/12$. The convergent local Bernoulli expansion gives
\begin{equation}\label{eq:hd-valuation}
 h_0(x)=\tfrac12+O(x),\qquad
 h_d(x)=\frac{B_{2d}}{(2d)!}x^{2d-1}+O(x^{2d+1})\quad(d\ge1).
\end{equation}
Write $P_d(x)=\sum_{j\in E_d}\alpha_{d,j}x^j$, where
$E_0=\{-1\}$ and
$E_d=\{-1,0,1,3,\ldots,2d-3\}$ for $d\ge1$, with the evident empty
positive portion for $d=1$.

\begin{definition}\label{def:Fd}
The Gamma-normalized compensated Hurwitz family is
\begin{equation}\label{eq:Fd-laplace}
 \mathcal F_d(s,z)=\int_0^\infty x^{s-1}e^{-zx}h_d(x)\dd x,
 \qquad\Rea s>-\nu_d,\quad\Rea z>0.
\end{equation}
Equivalently, on this domain it is the combined continuation of
\begin{equation}\label{eq:Fd-closed}
 \mathcal F_d(s,z)=\Gamma(s)\zeta(s,z)
       -\sum_{j\in E_d}\alpha_{d,j}\Gamma(s+j)z^{-s-j}.
\end{equation}
We also set $\mathcal R_d(s,z)=\mathcal F_d(s,z)/\Gamma(s)$ by
continuation. In particular,
\begin{align}
 \mathcal R_0(s,z)&=Z_s(z),\notag\\
 \mathcal R_d(s,z)&=Z_s(z)-\tfrac12z^{-s}
  -\sum_{k=1}^{d-1}\frac{B_{2k}}{(2k)!}(s)_{2k-1}z^{1-s-2k}
          \quad(d\ge1).\label{eq:Rd}
\end{align}
\end{definition}

\begin{proposition}[Holomorphy, zeros, and the position ladder]
\label{thm:Fd-domain}
The family $\mathcal F_d$ is jointly holomorphic on the domain in
\eqref{eq:Fd-laplace}. It is a holomorphic $L^2$ vector on the vertical
line $z=a+\ii y$ whenever
\begin{equation}\label{eq:Fd-L2}
 \Rea s>\frac12-\nu_d,\qquad\Rea a>0.
\end{equation}
It satisfies
\begin{equation}\label{eq:Fd-ladder}
 \partial_z\mathcal F_d(s,z)=-\mathcal F_d(s+1,z).
\end{equation}
For every integer $n$ with $0\le n<\nu_d$,
\begin{equation}\label{eq:Rd-zero}
 \mathcal R_d(-n,z)=0,\qquad
 \mathcal F_d(-n,z)=\frac{(-1)^n}{n!}
                    \left.\partial_s\mathcal R_d(s,z)\right|_{s=-n}.
\end{equation}
\end{proposition}
\begin{proof}
Near zero the kernel in \eqref{eq:Fd-laplace} is
$O(x^{\Rea s+\nu_d-1})$, and at infinity the exponential dominates
all polynomial terms in $P_d$. These bounds, with extra powers of
$\log x$ on compact spectral subsets, prove holomorphy. Squaring the
near-zero kernel gives exactly the sufficient $L^2$ condition
\eqref{eq:Fd-L2}. Plancherel then transfers the vector-valued holomorphy
to the vertical line. Multiplication by $-x$ proves the position ladder.

Formula~\eqref{eq:Fd-closed} follows by evaluating the separate terms
first for sufficiently large $\Rea s$, then continuing.
The function in \eqref{eq:Rd} is entire in $s$, after the original
$s=1$ pole cancellation. Since $1/\Gamma(s)$ has a simple zero at
$s=-n$ while $\mathcal F_d$ is holomorphic there, $\mathcal R_d(-n,z)=0$.
Finally,
$\Gamma(-n+\epsilon)=(-1)^n/(n!\epsilon)+O(1)$ gives the second formula.
\end{proof}

The kernel $L^2$ threshold is exact for an individual transform in its
Laplace domain, because the leading coefficients in
\eqref{eq:hd-valuation} are nonzero. We do not claim that the rectangular
pair of such thresholds is the maximal possible region for every product.
The theorem deliberately states a uniform domain in which all orders of
parameter differentiation are justified.

\subsection{Finite closure at every compensation depth}
Define
\begin{equation}\label{eq:J-integral}
 \mathcal J_{d,e}(w,A)=\int_0^\infty
          x^{w-1}e^{-Ax}h_d(x)h_e(x)\dd x.
\end{equation}
It is holomorphic for
$\Rea w>-\nu_d-\nu_e$, $\Rea A>0$.

\begin{theorem}[Arbitrary-depth product closure]\label{thm:depth-closure}
For $\Rea s>1/2-\nu_d$, $\Rea t>1/2-\nu_e$, and
$\Rea a,\Rea b>0$,
\begin{equation}\label{eq:Fd-pair}
 \frac1{2\pi}\int_{\R}\mathcal F_d(s,a+\ii y)
                         \mathcal F_e(t,b-\ii y)\dd y
       =\mathcal J_{d,e}(s+t-1,a+b).
\end{equation}
The right-hand side has the finite expression
\begin{align}
 \mathcal J_{d,e}(w,A)={}&
 \Gamma(w)\{\zeta(w-1,A)+(1-A)\zeta(w,A)\}\notag\\
 &-\sum_{j\in E_d}\alpha_{d,j}\Gamma(w+j)\zeta(w+j,A)\notag\\
 &-\sum_{k\in E_e}\alpha_{e,k}\Gamma(w+k)\zeta(w+k,A)\notag\\
 &+\sum_{j\in E_d}\sum_{k\in E_e}
   \alpha_{d,j}\alpha_{e,k}\Gamma(w+j+k)A^{-w-j-k}.
                  \label{eq:J-finite}
\end{align}
Every apparent pole of this expression in
$\Rea w>-\nu_d-\nu_e$ is removable. The identity and all its spectral
and position derivatives are identities of ordinary convergent pairings.
\end{theorem}
\begin{proof}
Equation~\eqref{eq:Fd-pair} is \cref{thm:plancherel} applied to
\eqref{eq:Fd-laplace}. For $\Rea w>2$, expand
\[
 (Q-P_d)(Q-P_e)=Q^2-P_dQ-P_eQ+P_dP_e.
\]
Every individual Mellin integral then converges. The $Q^2$ term was
evaluated earlier; the monomial--$Q$ terms yield
$\Gamma(w+j)\zeta(w+j,A)$, and the monomial products yield the last
line of \eqref{eq:J-finite}. The original combined integrand vanishes
to order $\nu_d+\nu_e$ at zero, proving its holomorphy on the larger
half-plane. Analytic continuation therefore cancels \emph{all} of the
apparent poles there, not merely those checked at finitely many indices.
The vector-valued argument proves the derivative assertion.
\end{proof}

For $d=e=0$, \eqref{eq:J-finite} is exactly
$\mathcal J_{0,0}(w,A)=\Gamma(w)K(w,A)$. Higher depths do not replace the
special-function alphabet: they enlarge its finite range of spectral
arguments while improving the admissible left half-plane. This is the
mechanism that permits antiderivatives to be paired without a divergent
vertical integral.

\section{Stirling remainders and normalized antiderivatives}
\label{sec:primitives}

\subsection{The first negative-order pairing}
At depth one, \eqref{eq:Rd-zero} and the standard zeta values at zero give
\begin{equation}\label{eq:Omega-F}
 \mathcal F_1(0,z)=\Omega(z)
    =\int_0^\infty e^{-zx}\frac{h_1(x)}x\dd x.
\end{equation}
For completeness, differentiate
$\mathcal R_1(s,z)=\zeta(s,z)-z^{1-s}/(s-1)-\frac12z^{-s}$ at $s=0$.
The result is
$\zeta'(0,z)-z\log z+z+\frac12\log z$, which is
\eqref{eq:Omega-intro}. Thus \eqref{eq:Omega-F} also follows directly
from the compensated Hurwitz family, independently of quoting Binet's
formula; it agrees with the classical representation in
\cite{dlmf-gamma}.

\begin{theorem}[Exact Stirling-remainder pairing]\label{thm:Omega-pair}
Let $\Rea a,\Rea b>0$ and $A=a+b$. Then
\begin{align}
 \frac1{2\pi}\int_{\R}\Omega(a+\ii y)\Omega(b-\ii y)\dd y
   ={}&A\zeta'(-1,A)-2\zeta'(-2,A)\notag\\
     &+\left(\frac{A^3}{6}-\frac{A^2}{2}+\frac A4\right)\log A
           +\frac{A^3}{36}+\frac A{12}.
                  \label{eq:Omega-pair}
\end{align}
\end{theorem}
\begin{proof}
The domain condition \eqref{eq:Fd-L2} holds at $d=1$, $s=0$.
Hence the left-hand side equals $\mathcal J_{1,1}(-1,A)$.
Since $h_1^2=h^2-h+1/4$, write
$\mathcal J_{1,1}(w,A)=\Gamma(w)\mathcal B(w,A)$, where
\begin{align}
 \mathcal B(w,A)={}&\frac{w-3}{w-1}\zeta(w-1,A)-A\zeta(w,A)
       +\frac{A^{2-w}}{(w-1)(w-2)}\notag\\
       &+\frac{A^{1-w}}{w-1}+\frac14A^{-w}.
                      \label{eq:B-Stirling}
\end{align}
At $w=-1$,
\[
 \mathcal B(-1,A)=2\zeta(-2,A)-A\zeta(-1,A)
                       +\frac{A^3}{6}-\frac{A^2}{2}+\frac A4=0.
\]
This is the Bernoulli-polynomial identity
$\zeta(-n,A)=-B_{n+1}(A)/(n+1)$ at $n=1,2$.
Because $\Res_{w=-1}\Gamma(w)=-1$,
\[
 \mathcal J_{1,1}(-1,A)=-\partial_w\mathcal B(-1,A).
\]
Differentiation before simplification gives
\begin{align*}
 -\partial_w\mathcal B(-1,A)={}&
 A\zeta'(-1,A)-2\zeta'(-2,A)-\tfrac12\zeta(-2,A)\\
 &+A^3\left(\frac{\log A}{6}-\frac5{36}\right)
       -\frac{A^2\log A}{2}+\frac{A^2}{4}+\frac{A\log A}{4}.
\end{align*}
Finally, $\zeta(-2,A)=-A^3/3+A^2/2-A/6$ yields
\eqref{eq:Omega-pair}.
\end{proof}

At $A=1$, the functional equation of the Riemann zeta function gives
$\zeta'(-2)=-\zeta(3)/(4\pi^2)$, proving
\eqref{eq:headline-stirling}. The sign in front of $\partial_w\mathcal B$
is forced by the negative residue of $\Gamma$ at $-1$.
Discarding the zero of $\mathcal B$ before differentiating would incorrectly
turn a nonzero norm into zero.

\subsection{A canonical primitive at every admissible depth}
Define the depth-$d$ Stirling remainder, for $d\ge1$, by
\begin{equation}\label{eq:Omega-depth}
 \Omega_d(z)=\Omega(z)-\sum_{k=1}^{d-1}
          \frac{B_{2k}}{2k(2k-1)}z^{1-2k}.
\end{equation}
The Gamma recurrence in \eqref{eq:Fd-closed} gives
$\mathcal F_d(0,z)=\Omega_d(z)$.

\begin{theorem}[Normalized primitive ladder and exact pairings]
\label{thm:primitive-ladder}
For $d\ge1$ and $0\le n<\nu_d=2d-1$, let
\begin{equation}\label{eq:primitive-definition}
 \mathcal P_{d,n}(z)=(-1)^n\mathcal F_d(-n,z)
      =\frac1{n!}\left.\partial_s\mathcal R_d(s,z)\right|_{s=-n}.
\end{equation}
Then
\begin{equation}\label{eq:primitive-ladder}
 \mathcal P_{d,0}=\Omega_d,\qquad
 \partial_z\mathcal P_{d,n}=\mathcal P_{d,n-1}\quad(n\ge1).
\end{equation}
Each $\mathcal P_{d,n}(x)$ tends to zero as real $x\to+\infty$.
This normalization uniquely specifies it among the holomorphic
$n$-fold antiderivatives of $\Omega_d$.
For $d,e\ge1$, $0\le n<\nu_d$, $0\le m<\nu_e$,
\begin{equation}\label{eq:primitive-pair}
 \frac1{2\pi}\int_{\R}\mathcal P_{d,n}(a+\ii y)
                       \mathcal P_{e,m}(b-\ii y)\dd y
       =(-1)^{n+m}\mathcal J_{d,e}(-n-m-1,a+b).
\end{equation}
The right-hand side is the removable value of the finite expression
\eqref{eq:J-finite}; no new integral is left unevaluated.
\end{theorem}
\begin{proof}
The definition and ladder follow from \cref{thm:Fd-domain}. The Laplace
representation is
\[
 \mathcal P_{d,n}(z)=(-1)^n\int_0^\infty x^{-n-1}e^{-zx}h_d(x)\dd x.
\]
Its near-zero power is $x^{\nu_d-n-1}$, with $\nu_d-n>0$.
Rescaling $x$ by a positive real $z$ and using an integrable split of the
kernel proves decay $O(z^{n-\nu_d})$ as $z\to+\infty$.
The difference of two $n$-fold antiderivatives is a polynomial of degree
less than $n$, and such a polynomial cannot tend to zero unless it is zero.
For $n=0$ the uniqueness statement simply identifies the function itself.

The integer restriction $n<\nu_d$ implies
$-n>1/2-\nu_d$; similarly for $m$. Thus \cref{thm:depth-closure} applies
and proves \eqref{eq:primitive-pair}. Moreover,
$-n-m-1>-\nu_d-\nu_e$, placing the spectral value strictly inside the
holomorphic domain of $\mathcal J_{d,e}$.
\end{proof}

The primitive formula is explicit in ordinary Hurwitz order derivatives:
\eqref{eq:Rd} involves only $\zeta(s,z)$, elementary powers, and a finite
Pochhammer polynomial. At $s=-n$ its first derivative therefore involves
$\zeta'(-n,z)$ and elementary terms. Pairing primitives at arbitrary
depths remains finite by \eqref{eq:J-finite}. At its integer resonances,
first derivatives of Hurwitz zeta and digamma values suffice; the next
section gives a stable finite-part rule. Higher \emph{spectral} jets may
require correspondingly higher zeta derivatives and Stieltjes constants,
but always only finitely many.

\section{Rationally related vertical scales}
\label{sec:rational-scales}

Unequal slopes introduce two arithmetic progressions into the product
kernel. For a rational slope ratio those progressions are commensurate,
and their apparent double-zeta structure reduces to a finite collection
of single Hurwitz zeta functions. The scale factors in this statement
are important: one cannot obtain it by replacing $a+b$ with a weighted
sum in \cref{thm:basic} while leaving all other factors unchanged.

Let $p,q$ be positive integers, not necessarily coprime. Put
\begin{equation}\label{eq:rational-grid}
 B=qa+pb,\qquad R=pq,\qquad
 c_{ij}=\frac{B+qi+pj}{R}\quad(0\le i<p,\ 0\le j<q),
\end{equation}
and define
\begin{equation}\label{eq:Dpq}
 D_{p,q}(w,B)=R^{-w}\sum_{i=0}^{p-1}\sum_{j=0}^{q-1}
       \{\zeta(w-1,c_{ij})+(1-c_{ij})\zeta(w,c_{ij})\}.
\end{equation}

\begin{lemma}[Finite residue-grid reduction]\label{thm:grid}
For $\Rea w>2$ and $\Rea B>0$,
\begin{equation}\label{eq:QQ-grid}
 \int_0^\infty x^{w-1}e^{-Bx}Q(qx)Q(px)\dd x
                 =\Gamma(w)D_{p,q}(w,B).
\end{equation}
\end{lemma}
\begin{proof}
Write each pair of nonnegative integers uniquely as
$n=p\ell+i$, $m=qk+j$ with $0\le i<p$, $0\le j<q$.
Then $qn+pm=qi+pj+pq(\ell+k)$, so
\[
 Q(qx)Q(px)=\sum_{i=0}^{p-1}\sum_{j=0}^{q-1}
              \frac{e^{-(qi+pj)x}}{(1-e^{-Rx})^2}.
\]
Expanding the square produces the multiplicity $\ell+k+1$.
Rescale $Rx$ and use the $Q^2$ calculation in \cref{thm:basic}.
Absolute convergence in the initial half-plane justifies every sum.
No coprimality condition is used in this decomposition.
\end{proof}

Define the compensated grid transform by
\begin{equation}\label{eq:Jpq-integral}
 \mathcal J_{d,e}^{p,q}(w,B)=\int_0^\infty
             x^{w-1}e^{-Bx}h_d(qx)h_e(px)\dd x.
\end{equation}

\begin{theorem}[Rational-scale closure at every depth]
\label{thm:rational-depth}
Under the same $L^2$ conditions on $s,t,a,b$ as in
\cref{thm:depth-closure},
\begin{align}
 &\frac1{2\pi}\int_{\R}\mathcal F_d(s,a+\ii py)
                              \mathcal F_e(t,b-\ii qy)\dd y\notag\\
 &\quad=q^{s-1}p^{t-1}\,
       \mathcal J_{d,e}^{p,q}(s+t-1,qa+pb).
                          \label{eq:rational-depth-pair}
\end{align}
The finite expression for the right-hand transform is
\begin{align}
 \mathcal J_{d,e}^{p,q}(w,B)={}&\Gamma(w)D_{p,q}(w,B)\notag\\
 &-\sum_{k\in E_e}\alpha_{e,k}p^k\Gamma(w+k)q^{-w-k}
                                 \zeta(w+k,B/q)\notag\\
 &-\sum_{j\in E_d}\alpha_{d,j}q^j\Gamma(w+j)p^{-w-j}
                                 \zeta(w+j,B/p)\notag\\
 &+\sum_{j\in E_d}\sum_{k\in E_e}
       \alpha_{d,j}\alpha_{e,k}q^jp^k\Gamma(w+j+k)B^{-w-j-k}.
                          \label{eq:Jpq-finite}
\end{align}
It is holomorphic for $\Rea w>-\nu_d-\nu_e$, $\Rea B>0$ after all
removable cancellations. The identities are stable under arbitrary
spectral and position derivatives inside the stated domains.
\end{theorem}
\begin{proof}
The scaled formula \eqref{eq:scaled-plancherel} pairs
$x^{s-1}e^{-ax}h_d(x)$ with $x^{t-1}e^{-bx}h_e(x)$.
Their values at $qx$ and $px$ contribute
$q^{s-1}p^{t-1}x^{s+t-2}e^{-(qa+pb)x}h_d(qx)h_e(px)$,
which proves \eqref{eq:rational-depth-pair}.
For $\Rea w>2$, expand
$(Q(qx)-P_d(qx))(Q(px)-P_e(px))$.
The first term is \cref{thm:grid}; a monomial from $P_e(px)$ contributes
$p^kx^k$ against $Q(qx)$ and hence the second line of
\eqref{eq:Jpq-finite}. The other two lines follow in the same way.
The combined kernel vanishes to order $\nu_d+\nu_e$ at zero.
Its holomorphy proves all pole cancellations and the derivative assertion
exactly as in \cref{thm:depth-closure}.
\end{proof}

If the actual slopes are $cp,cq$ with $c>0$, one multiplies the result
by $1/c$ after changing $y$ to $cy$. Thus the theorem covers every positive
rational slope ratio. At $p=q=c$ it is consistent with the ordinary
$1/c$ scaling law. For an irrational ratio the Laplace identity survives,
but the finite residue grid used in \cref{thm:grid} is unavailable; no
non-reducibility claim is implied.

\subsection{An explicit all-order rational-scale Stieltjes formula}
For depth zero write
$\mathcal J_{0,0}^{p,q}(w,B)=\Gamma(w)K_{p,q}(w,B)$.
Then \eqref{eq:Jpq-finite} simplifies to
\begin{align}
 K_{p,q}(w,B)={}&D_{p,q}(w,B)
 -\frac{p^{-1}q^{1-w}\zeta(w-1,B/q)
       +q^{-1}p^{1-w}\zeta(w-1,B/p)}{w-1}\notag\\
 &+\frac{B^{2-w}}{pq(w-1)(w-2)}.
                                  \label{eq:Kpq}
\end{align}
For integers $m,n\ge0$ define
\[
 I_{m,n}^{p,q}(a,b)=\frac1{2\pi}\int_{\R}
               G_m(a+\ii py)G_n(b-\ii qy)\dd y.
\]
Then the theorem gives
\begin{equation}\label{eq:rational-Stieltjes-jets}
 \boxed{\displaystyle
 I_{m,n}^{p,q}(a,b)=(-1)^{m+n}m!n!\,[u^mv^n]
       q^up^v C(u,v)K_{p,q}(1+u+v,B).}
\end{equation}
The exponential slope factors are essential even for first-order jets.

For an entirely pole-free coefficient recipe put
\begin{equation}\label{eq:coefficient-symbols}
 z_j(c)=\frac{\zeta^{(j)}(0,c)}{j!},\qquad
 g_j(c)=\frac{(-1)^j\gamma_j(c)}{j!},\qquad
 e_j(\lambda)=\frac{(-\log\lambda)^j}{j!}.
\end{equation}
Let
$K_{p,q}(1+\epsilon,B)=\sum_{k\ge0}\kappa_k^{p,q}(B)\epsilon^k$.

\begin{corollary}[Finite coefficient formula]\label{thm:rational-coeff}
For each integer $k\ge0$,
\begin{align}
 \kappa_k^{p,q}(B)={}&\frac1R\sum_{i=0}^{p-1}\sum_{j=0}^{q-1}
  \left\{\sum_{\ell=0}^{k}e_\ell(R)
       [z_{k-\ell}(c_{ij})+(1-c_{ij})g_{k-\ell}(c_{ij})]
      +(1-c_{ij})e_{k+1}(R)\right\}\notag\\
 &-\frac1p\sum_{\ell=0}^{k+1}e_\ell(q)z_{k+1-\ell}(B/q)
  -\frac1q\sum_{\ell=0}^{k+1}e_\ell(p)z_{k+1-\ell}(B/p)\notag\\
 &-\frac BR\sum_{\ell=0}^{k+1}e_\ell(B).
                                  \label{eq:rational-coeff}
\end{align}
Together with \eqref{eq:rational-Stieltjes-jets}, this evaluates every
rational-scale Stieltjes pairing by finitely many one-variable Hurwitz
jets, Stieltjes constants, elementary logarithms, and universal zeta
constants from $C$.
\end{corollary}
\begin{proof}
Expand \eqref{eq:Kpq} at $w=1+\epsilon$.
The first double sum contains a simple zeta pole multiplied by
$R^{-1}e^{-\epsilon\log R}$; its contribution to the coefficient of
$\epsilon^k$ is exactly $(1-c_{ij})e_{k+1}(R)/R$ in addition to the
regular convolution of Taylor coefficients. The two crossed terms
contain $e^{-\epsilon\log q}\zeta(\epsilon,B/q)/\epsilon$ and its
$p$ analogue. The last term is
$-(B/R)e^{-\epsilon\log B}/\{\epsilon(1-\epsilon)\}$.
Their coefficients are the last two lines of \eqref{eq:rational-coeff}.

To check the cancellation preceding extraction, observe that
\[
 \sum_{i,j}(1-c_{ij})=\frac{p+q}{2}-B.
\]
The total residue of \eqref{eq:Kpq} is therefore
\[
 \frac{(p+q)/2-B}{pq}
 -\frac1p\Bigl(\frac12-\frac Bq\Bigr)
 -\frac1q\Bigl(\frac12-\frac Bp\Bigr)-\frac B{pq}=0.
\]
This also checks the powers of $p,q$ in the coefficient formula.
\end{proof}

For $p=q=1$, \eqref{eq:rational-coeff} reduces exactly to
\eqref{eq:k-coefficients}. At general $p,q$, Stieltjes jets at the grid
points enter only up to order $m+n$, while zeta order derivatives at
$B/p$ and $B/q$ enter only up to order $m+n+1$.

\subsection{The two-scale quarter-Gamma identity}
The case $k=0$ of \eqref{eq:rational-coeff} is
\begin{align}
 K_{p,q}(1,B)={}&\frac1R\sum_{i,j}
   \{\zeta(0,c_{ij})-(1-c_{ij})\psi(c_{ij})
                              -(1-c_{ij})\log R\}\notag\\
 &-\frac1p\zeta'(0,B/q)-\frac1q\zeta'(0,B/p)
   +\frac{\log q}{p}\zeta(0,B/q)
   +\frac{\log p}{q}\zeta(0,B/p)\notag\\
 &+\frac BR(\log B-1).
                         \label{eq:Kpq-one}
\end{align}
For $p=1,q=2,a=b=1/2$, the grid is $3/4,5/4$ and $B=3/2$.
Using
$\psi(5/4)-\psi(3/4)=4-\pi$, the first line reduces to $-\pi/8$.
Thus
\begin{align*}
 K_{1,2}(1,3/2)={}&-\frac\pi8-\zeta'(0,3/4)
       -\frac12\zeta'(0,3/2)-\frac14\log2
       +\frac34\{\log(3/2)-1\}.
\end{align*}
Insert Lerch's formula for $\zeta'(0,z)$,
$\Gamma(3/2)=\sqrt\pi/2$, and
$\Gamma(1/4)\Gamma(3/4)=\pi\sqrt2$.
The result is \eqref{eq:headline-quarter}, proving that identity without
numerical recognition of its constants.

\section{Evaluating removable integer values without numerical cancellation}
\label{sec:finite-part-engine}

The formulas for deep compensation contain many singular summands even
when their sum is regular. A stable evaluation should compute their
Laurent constant terms, retaining derivatives of analytic prefactors.
The following rules are enough for every integer specialization of
\eqref{eq:J-finite} and \eqref{eq:Jpq-finite}.

Let $H_0=0$ and $H_n=\sum_{j=1}^n1/j$. For $n\ge0$,
\begin{align}
 \Gamma(-n+\epsilon)
  &=\frac{(-1)^n}{n!}
     \left\{\frac1\epsilon+H_n-\gamma+O(\epsilon)\right\},
                                 \label{eq:Gamma-laurent}\\
 \CT_{\epsilon=0}\bigl[\Gamma(-n+\epsilon)
                    \zeta(-n+\delta+\epsilon,C)\bigr]
  &=\frac{(-1)^n}{n!}
    \{\zeta'(-n+\delta,C)+(H_n-\gamma)\zeta(-n+\delta,C)\},
                                 \label{eq:GZ-negative-ct}
\end{align}
provided the zeta factor is regular there. The needed unshifted grid
terms have $\delta=0,-1$, so this proviso always holds when their Gamma
factor has a nonpositive integer argument. If a regular Gamma factor
meets the zeta pole instead, $k\ge1$ and $k+\delta=1$, then
\begin{equation}\label{eq:GZ-positive-ct}
 \CT_{\epsilon=0}[\Gamma(k+\epsilon)\zeta(1+\epsilon,C)]
       =\Gamma(k)\{\psi(k)-\psi(C)\},
\end{equation}
with residue $\Gamma(k)$. The elementary terms have
\begin{equation}\label{eq:GE-ct}
 \CT_{\epsilon=0}[\Gamma(-n+\epsilon)C^{n-\epsilon}]
       =\frac{(-1)^nC^n}{n!}\{H_n-\gamma-\log C\}.
\end{equation}
All these formulas follow by multiplication of one-term Laurent and
Taylor expansions.

The prefactor rule is indispensable. If a singular factor has residue
$r$ and constant term $c$, and
$\theta(\epsilon)=\theta_0+\theta_1\epsilon+O(\epsilon^2)$, then
\begin{equation}\label{eq:prefactor-rule}
 \CT\left[\theta(\epsilon)\left(\frac r\epsilon+c+O(\epsilon)\right)\right]
        =\theta_0c+\theta_1r.
\end{equation}
For example, $q^{-w}$ contributes the logarithmic derivative $-\log q$.
Dropping the second term in \eqref{eq:prefactor-rule} would erase the
very slope logarithms required by \eqref{eq:rational-coeff}.

\subsection{Finite algorithm and its justification}
To evaluate $\mathcal J_{d,e}^{p,q}(w_0,B)$ at an integer $w_0$ inside its
holomorphic domain, expand each summand of \eqref{eq:Jpq-finite} only
through its residue and constant term. For grid summands, use
\eqref{eq:GZ-negative-ct} or \eqref{eq:GZ-positive-ct}, and include the
factor $R^{-w}$ by \eqref{eq:prefactor-rule}. For crossed terms, shift
$w_0$ by the integer monomial exponent and include the corresponding
$q^{-w}$ or $p^{-w}$ factor. For elementary terms, use
\eqref{eq:GE-ct}. Sum the resulting residues and constant terms.
\cref{thm:rational-depth} proves the residue sum is exactly zero and
that the constant sum is the value of the ordinary convergent integral.
Thus a finite part is used solely as a computational operation on a
meromorphic formula whose total singularity is removable.

The function \texttt{J\_at\_integer} in the accompanying code implements
this procedure, returning both the finite value and the residual sum.
It avoids taking tiny spectral offsets and subtracting enormous nearly
equal quantities. First-order Hurwitz order derivatives suffice for these
integer values; higher spectral derivatives are obtained by retaining
more Laurent coefficients. This gives an implementable extension of the
primitive calculus rather than leaving its resonances as formal limits.

\section{Audit findings, corrections, and verification status}
\label{sec:audit}

\subsection{Scope of the source review}
The inspected canonical convolution foundation in \cite{repo} works with
the mean-zero periodic family having nonzero Fourier coefficients
$(2\pi\ii k)^{-\alpha}$. Its contact and finite-part conventions belong
to that geometry. They have not been transferred silently to the real
vertical line. In the present setting, the elementary term subtracted from
$\zeta(s,z)$ is $z^{1-s}/(s-1)$, and the analytic mechanism is a Laplace
kernel on $(0,\infty)$.

No new false theorem was established in the canonical source portions
actually inspected. It would therefore be inappropriate to issue a
blanket erratum against the existing manuscript. The practical correction
found during this continuation concerns the complex-argument numerical
backend described next. We also isolate three normalization mistakes
that would invalidate an attempted extension of the earlier calculus:
\begin{enumerate}
\item Subtracting only the spectral pole $1/(s-1)$ does not remove spatial
      growth. At $s=1$ it leaves $\gamma_0(z)=-\psi(z)$; the squared
      vertical norm diverges because $\psi(a+\ii y)$ has logarithmic
      growth. Adding $\log z$ produces the convergent $G_0$ pairing.
\item The zero $\mathcal R_d(-n,z)=0$ does not imply that its spectral
      derivative vanishes. The pole of $\Gamma(s)$ turns that derivative
      into the nonzero value in \eqref{eq:Rd-zero}. The Stirling norm
      illustrates the corresponding zero-times-pole cancellation in
      the paired expression.
\item Replacing analytic $\log\Gamma(z)$ by the principal logarithm of
      $\Gamma(z)$ changes its branch by discontinuous multiples of
      $2\pi\ii$. It destroys the Laplace representation and need not
      preserve any vertical norm in this paper.
\end{enumerate}
These are explicit scope guards, not allegations that the inspected
canonical chapters made these mistakes.

\subsection{A reproduced mpmath 1.3.0 nonreal-argument failure}
\label{sec:software-audit}
During the independent numerical replay, the installed
\texttt{mpmath 1.3.0} function \texttt{stieltjes(0,z)} failed the exact
calibration identity $\gamma_0(z)=-\psi(z)$ at $z=0.8+0.6\ii$.
At 60 decimal digits of working precision it returned approximately
\[
 0.4662545368664062120-0.9931920114430890458\ii,
\]
whereas $-\psi(z)$ is approximately
\[
 0.4278547847547755586-1.0184230001203705136\ii.
\]
The absolute discrepancy was approximately
$0.04594718437366742995$, not a last-digit effect.
The version-tagged source \cite{mpmath13} takes a real part inside its
integrand. That real-parameter formula is not a valid analytic continuation
in a nonreal argument.

The repair used here is explicit. Complex Stieltjes values are checked by
the holomorphic paired Hermite formula \eqref{eq:complex-Hermite} and by
Cauchy coefficients of the entire function $u\mapsto Z_{1+u}(z)$.
Positive-real arguments in \eqref{eq:k-coefficients} may still use the
installed Stieltjes routine; an independent positive-real control was
included. The code does not claim this finding is new, that later mpmath
versions share it, or that any particular canonical repository result
was affected. A corpus-wide search for such unsafe nonreal calls would
be a separate audit task.

For the Cauchy check, with $0<\rho<1$ and integer $N>m$, the approximation
used is
\begin{equation}\label{eq:Cauchy-diagnostic}
 G_m(z)\approx\frac{(-1)^m m!}{N\rho^m}
   \sum_{k=0}^{N-1}Z_{1+\rho\omega_k}(z)\omega_k^{-m},
 \qquad \omega_k=e^{2\pi\ii k/N}.
\end{equation}
The actual replay used $\rho=1/8$ and $N=96$ at 60-digit precision for
orders zero through two. This is a trapezoidal discretization of Cauchy's
coefficient formula, not a certified error enclosure. Its agreement with
the independent Laplace and Hermite evaluations is recorded separately
from the analytic proof.

\subsection{What the verification establishes}
The supplied exact checks contain \textbf{538 passing assertions}:
7 kernel-valuation certificates, 36 compensated spectral-zero identities,
210 equal-scale residue cancellations, 234 rational-scale residue
cancellations, 49 finite-lattice polynomial identities, and 2
Stirling-polynomial reductions. They use exact rational arithmetic and
symbolic Bernoulli polynomials. These finite checks would detect many
sign, indexing, or scale errors, but the proofs above, not the finite
range of tests, establish the all-order assertions.

The numerical replay contains \textbf{76 passing comparisons} at
\textbf{60 decimal digits of working precision}, using acceptance
criterion
\[
 \frac{|\text{left}-\text{right}|}{\max(1,|\text{left}|,|\text{right}|)}
      <10^{-35}.
\]
The largest recorded scaled discrepancy was approximately
$2.04\times10^{-53}$, in a basic spectral moment with $w=0.8$, $A=1$.
The test set includes direct vertical-line quadrature, independent
Laplace moments, mixed Stieltjes jets, complex Cauchy coefficients,
polygamma integrals, rational scales, negative-order compensated
moments, removable integer values, primitive derivatives, Lerch
confluence, root-of-unity Gamma reductions, and finite harmonic sums.
Four of these comparisons form the separate complex-safe replay; its
negative control reproduces the dependency limitation just discussed.

These are unvalidated floating-point comparisons, not interval proofs.
For direct vertical integrals, the implementation switches to the
standard asymptotic expansions of $G_0$ and $\Omega$ beyond $|z|=60$
to avoid catastrophic subtraction at very large quadrature nodes.
The resulting quadrature remains a diagnostic. No Lean or other
proof-assistant formalization is supplied.

\subsection{Reproduction and proposed integration}
The ZIP contains the modular article sources, a standalone TeX source,
the compiled PDF, exact and numerical scripts, machine-readable results,
a version-specific negative-control record, source audit, proof-status
ledger, build instructions, and checksums. The core commands are
\begin{verbatim}
python verification/exact_checks.py
python verification/numerical_checks.py core
python verification/numerical_checks.py extensions
python verification/complex_safe_replay.py
latexmk -pdf -interaction=nonstopmode -halt-on-error article.tex
\end{verbatim}

A suitable integration is a new chapter or research note on
\emph{vertical-line Hurwitz pairings}, placed near the existing integration
and log-Gamma material, not merged into the periodic contact-term chapter
as if its integrals had the same meaning. Suggested labels and dependencies
are in \texttt{INTEGRATION\_NOTES.md}. The incoming ZIP should be retained
unchanged as a provenance artifact, consistent with the repository's
incoming-report policy. No repository write was performed in preparing
this package.

\section{Further research questions}
\label{sec:questions}

The proved results suggest the following concrete extensions. These are
questions, not additional theorems claimed by the present paper.

\begin{question}[Irrational-scale closure and its precise obstruction]
For $p/q\notin\mathbb Q$, classify when
$\int_0^\infty x^{w-1}e^{-Bx}h_d(qx)h_e(px)\dd x$ admits a finite
representation by terms of the form
$\Gamma(w+j)\lambda^{-w}\zeta(w+j,C)$ and elementary Gamma--power terms,
with parameters independent of $w$. The allowed coefficient field and
allowed scales must be fixed before asserting a non-reduction theorem.
The Laplace identity is already available; what is missing is a rigorous
classification of the corresponding exponential-kernel algebra.
\end{question}

\begin{question}[Joint color coalescence and Bernoulli compensation]
Develop a compensation hierarchy for
$(1-\xi e^{-x})^{-1}$ that remains uniform as $\xi\to1$ and produces
well-defined mixed color/spectral jets. Can its exact transition terms
be written in the same finite Hurwitz alphabet as
\eqref{eq:rational-coeff}? A formula should specify the order of color
coalescence and spectral differentiation rather than rely on formal
substitution into a singular divided difference.
\end{question}

\begin{question}[Higher products and multiple-zeta closure]
For three or more factors with both vertical orientations, Plancherel
produces convolution constraints rather than a single diagonal product.
Determine the smallest class of multiple Hurwitz or Tornheim-type
functions closed under the resulting compensated integrals, and identify
integer resonances where that class collapses to ordinary Hurwitz order
jets. The two-factor result here provides exact normalization tests for
such a theory.
\end{question}

\begin{question}[Minimal arithmetic bases for rational scales]
The $pq$-point grid is explicit but not always minimal. Use Gamma
multiplication, reflection, and finite character decompositions to find
minimal provable spanning sets for
$I_{m,n}^{p,q}(a,b)$ at rational $a,b$. The quarter-Gamma formula
\eqref{eq:headline-quarter} is the first simple instance. Any assertion
of a basis must distinguish a proved spanning statement from unproved
linear independence of its constants.
\end{question}

\begin{question}[Exact reductions of remaining integral conjectures]
Determine which still-conjectural polylogarithm integrals in the broader
repository can be transformed into the finite-subtraction kernel algebra
of this paper. A successful reduction would need an explicit transform
identity and branch/convergence proof. No such reduction of $S_6$ or the
revised $S_8$ candidate is supplied here, and numerical agreement alone
would not constitute one.
\end{question}

\begin{question}[Certified complex jets and formalization]
Replace the floating-point Cauchy replay with an interval-certified
complex-zeta coefficient engine, including explicit contour truncation
and rounding errors. In parallel, formalize the bilinear
Laplace--Plancherel lemma, compact-parameter domination, and finite kernel
identity. These three components would turn the present proof ledger
into a reusable formal route to arbitrary-depth and rational-scale
Stieltjes identities.
\end{question}

\paragraph{Conclusion.}
The central mathematical outcome is finite exact closure for a genuinely
different integration geometry from the existing periodic theory. Spatial
compensation gives all Stieltjes indices; deeper Bernoulli compensation
gives Gamma remainders and normalized primitives; commensurate vertical
scales give finite arithmetic grids; and colored kernels give explicit
polylogarithm order identities. The difficult bookkeeping points---joint
parameter holomorphy, removable poles, slope logarithms, and complex
branches---are part of the proofs, not delegated to a numerical fit.

\clearpage
\begin{thebibliography}{99}
\bibitem{repo}
V. Reshetnikov, \emph{ProveIt}, public research repository, canonical
polylogarithm manuscript and incoming-report directory. Inspected source
paths and blob identifiers are recorded in \texttt{SOURCE\_AUDIT.md}.
\url{https://github.com/VladimirReshetnikov/ProveIt}.
\bibitem{dlmf-zeta}
NIST Digital Library of Mathematical Functions, \S25.11, \emph{Hurwitz Zeta
Function}, especially the integral representation, special values,
and parameter-differentiation formulas.
\url{https://dlmf.nist.gov/25.11}.
\bibitem{dlmf-gamma}
NIST Digital Library of Mathematical Functions, \S5.9, \emph{Integral
Representations}, including Binet's formula and digamma integrals.
\url{https://dlmf.nist.gov/5.9}.
\bibitem{dlmf-lerch}
NIST Digital Library of Mathematical Functions, \S25.14, \emph{Lerch's
Transcendent}.
\url{https://dlmf.nist.gov/25.14}.
\bibitem{shpot2016}
M. A. Shpot, M. P. Chaudhary, and R. B. Paris,
\emph{Integrals of products of Hurwitz zeta functions and the Casimir effect
in $\phi^4$ field theories}, arXiv:1603.00722 (2016).
\url{https://arxiv.org/abs/1603.00722}.
\bibitem{shpot-paris}
M. A. Shpot and R. B. Paris,
\emph{Integrals of products of Hurwitz zeta functions via Feynman
parametrization and two double sums of Riemann zeta functions},
arXiv:1609.05658 (2016).
\url{https://arxiv.org/abs/1609.05658}.
\bibitem{mpmath13}
The mpmath developers, \emph{mpmath}, version 1.3.0,
\texttt{mpmath/functions/zeta.py}, function \texttt{stieltjes}.
Version-specific source inspected for the nonreal-argument diagnostic.
\url{https://github.com/mpmath/mpmath/blob/1.3.0/mpmath/functions/zeta.py}.
\end{thebibliography}
\end{document}
